\section{Introduction}

\label{sec:introduction}

\subsection{Motivation}

\label{sec:introduction-motivation}

Tool-using language-model agents turn generated text into actions. An indirect prompt injection embedded in a message, document, webpage, memory item, or tool response can therefore become a request to send data, modify state, or invoke a privileged operation. Public environments such as AgentDojo and InjecAgent show that security must be evaluated together with task utility rather than as a text-classification problem alone \citep{L09,L10}. Runtime systems including VIGIL, DRIFT, the Task Shield, IPIGuard, and AgentSpec consequently interpose verification, task alignment, dependency, information-flow, or policy checks before an effect is committed \citep{L01,L02,L03,L20,L25}.

This literature also makes a broad novelty claim untenable. Verify-before-commit, contextual policy generation, least privilege, planned dependency checks, and control/data-flow separation are already occupied design spaces \citep{L01,L08,L24,L25,L26}. Recent preprints go further by auditing context influence, comparing an authorization graph with a parameter-source execution graph, or controlling source authority in next-action decisions \citep{L21,L22,L35}. Evidence-backed permission graphs with formal runtime authorization further narrow the mechanism claim boundary \citep{L36}. The open question for this study is therefore not whether a provenance-aware guard can be built for the first time. It is whether authorization information, enforcement structure, and failure handling materially change the security and utility measured for tool-using agents.

\subsection{Evaluation gap}

\label{sec:introduction-evaluation-gap}

Three distinctions are easy to blur in an agent-security experiment.

First, a detector score or model refusal is not an executed-effect endpoint. Prompt-level structuring, spotlighting, instruction hierarchy, and injection detection can reduce model compliance, but they do not by themselves prove that an un\-au\-thor\-ized action failed to cross an external runtime boundary \citep{L07,L30,L31,L32}.

Second, the authorization regime is part of the treatment. A guard that sees per-sample answer keys, evaluator labels, or an exact policy reconstructed from the expected answer is solving a different problem from a guard that receives only trusted user intent and a non-oracle policy compiled before execution. The present source audit does not establish that such leakage is universal, so we frame this as authorization uncertainty rather than a domain-wide answer-key defect.

Third, low attack success can be purchased by blocking useful work. A runtime defense must therefore report executed unsafe effects, benign task completion, system confirmation-request incidence, operational failures, latency, tokens, and cost under a common intention-to-treat denominator. Static attacks are also insufficient for a robustness claim when an attacker can adapt to the public defense interface \citep{L27}.

\subsection{Research questions}

\label{sec:introduction-research-questions}

The study asks six questions.

\begin{enumerate}

\item \textbf{RQ1: Authorization regime.} How do \nolinkurl{exact_gold}, \nolinkurl{auto_compiled}, \nolinkurl{broad}, \nolinkurl{underspecified}, and \nolinkurl{conflicting} authorization regimes change unsafe effect, benign task completion, and \nolinkurl{ASK/DENY} burden?

\item \textbf{RQ2: Prompt and runtime components.} What are the independent main effects of the prompt profile and the runtime module bundle, and do they interact?

\item \textbf{RQ3: Conditional robustness and model replication.} Are the factorial effect directions consistent in a representative sample, an outcome-enriched challenge set, and a same-provider Flash replication?

\item \textbf{RQ4: Source authority.} When proposition, value, position, and policy are held fixed, how does moving a critical argument between trusted and untrusted origins change model--Guard behavior?

\item \textbf{RQ5: Adaptive tail risk.} Among attacks with no unsafe effect by budget 200, how often does a first unsafe effect appear by budget 500?

\item \textbf{RQ6: Contemporaneous attack activity.} In a historically risk-enriched cohort, do control attacks remain active in the current execution block, and how do six runtime conditions compare under the same provider window?

\end{enumerate}

\subsection{Study design}

\label{sec:introduction-study-design}

M3 uses a three-state runtime decision (\nolinkurl{ALLOW}, \nolinkurl{ASK}, \nolinkurl{DENY}), a projection that excludes evaluator-only fields from online state, permission monotonicity for automatic paths, parameter-source evidence, and independent confirmation for privilege expansion. The completed authorization experiment crosses a policy package and a full package with non-oracle \nolinkurl{auto_compiled} and oracle-sensitivity \nolinkurl{exact_gold}; a nested subset adds three authorization stress regimes. The incremental factorial experiment then holds \nolinkurl{auto_compiled} fixed and crosses prompt profile (\nolinkurl{normal}/\nolinkurl{safe}) with runtime modules (\nolinkurl{policy}/\nolinkurl{full}). A source-authority counterfactual holds text, parameter value, relative position, and policy fixed while changing only provenance authority. Separate challenge, Flash-replication, adaptive-tail, and single-block activity calibration studies test conditional robustness. Timeouts, parsing failures, refusals, maximum-step outcomes, and provider failures remain in the intention-to-treat denominator.

The authorization-regime sample contains 180 paired cases and 900 DeepSeek V4 Pro work items. The incremental design adds 720 Pro factorial runs, 240 Pro source-authority runs plus 480 deterministic Guard decisions, 240 Pro challenge runs, 360 Flash replication runs, and a conditional adaptive tail of at most 3,600 additional queries. A final calibration adds 360 DeepSeek work items over 60 historically control-positive attacks and six concurrent conditions. The design separately reuses the completed 480-record medium study, 5,082-record Pro Public matrix, 1,791-record Flash Public matrix, and the 40-task adaptive evidence at budgets 50 and 200. Two bounded Gemini attempts remain diagnostic only. Provider tools and real external effects are prohibited throughout.

\subsection{Empirical motivation and study framing}

\label{sec:introduction-empirical-motivation-and-study-framing}

The completed medium falsification study did not support full-package superiority. The full package reduced unsafe effects from 8/36 under the structural condition to 0/36, but the policy package also produced 0/36. Benign tool completion fell from 12/24 under the structural condition to 6/24 under the full package, violating the preregistered non-inferiority margin. These observations motivate an empirical study of authorization uncertainty and the security--utility frontier rather than a claim that the full package is a new or state-of-the-art defense.

The larger Public primary matrix reinforced this finding. Across 600 attacks, the full package reduced unsafe effects from 133 under the structural condition to 38, but the policy package produced only 13. Across 168 benign tool tasks, both M3 variants completed 49 versus 103 under the structural condition. The full package therefore shows neither incremental security nor utility value beyond the policy package in the primary-model evidence. The 1.4955\% retained diagnostic fraction also exceeds the predefined 1\% threshold for a clean configuration-quality claim.

The completed 40-task adaptive panel observed first unsafe effect for 4/40 tasks by budget 50 and 18/40 by budget 200. These points show substantial budget sensitivity but do not establish saturation. The incremental tail extended only a prespecified subset of 12 B200 survivors toward 500. No new first success was observed (0/12; Wilson 95\% upper limit 24.25\%); this is a conditional analysis rather than a complete B500 curve or saturation result.

The prospectively frozen authorization-regime experiment is now complete for DeepSeek. Unsafe effects remained at or near zero across all four main cells, but exact authorization produced opposite benign-tool directions across the two benchmarks and increased failure-sensitive attack risk on AgentDojo. It also replaced ASK with more frequent DENY and increased latency. The full package showed no stable advantage over the policy package. These unfavorable results are retained. A separate risk-enriched single-block calibration found 17/60 unsafe effects under no defense, confirming contemporaneous activity in the selected cohort; full/safe produced 0/60, but failures and the enriched sampling frame preclude deployment-prevalence, utility, or temporal-repeatability claims. Gemini free-tier execution remained too incomplete in both attempted studies for a cross-provider effect estimate.

\subsection{Contributions}

\label{sec:introduction-contributions}

The current evidence supports six bounded contributions:

\begin{enumerate}

\item an evaluator/online-state contract that makes authorization visibility explicit and prevents evaluator-only fields from entering runtime state;

\item a contemporaneous factorial evaluation that identifies prompt-profile and runtime-module-bundle effects separately;

\item a controlled source-authority counterfactual that separates deterministic Guard decisions from end-to-end action and utility outcomes;

\item an intention-to-treat experiment that records official model identity, request commitments, isolated effects, and terminal diagnostics while retrying only response-less failed rows; and

\item an empirical account of authorization regime, confirmation/blocking burden, outcome-enriched instability, same-provider replication, and conditional defense-aware attack budget; and

\item a contemporaneous activity calibration that distinguishes an active control attack set from low event rates under defended conditions without treating a historical risk-enriched sample as deployment prevalence.

\end{enumerate}

These contributions do not assert first, only, or SOTA status. Preprints are used to define the contemporary novelty boundary, not as peer-reviewed evidence.
